\documentclass[11pt]{article}

\usepackage[margin=1in]{geometry}
\usepackage{amsmath,amssymb,amsthm,mathtools}
\usepackage{booktabs}
\usepackage{needspace}
\usepackage[hidelinks,hypertexnames=false]{hyperref}
\usepackage{microtype}

\newtheorem{theorem}{Theorem}[section]
\newtheorem{corollary}[theorem]{Corollary}
\newtheorem{proposition}[theorem]{Proposition}
\newtheorem{lemma}[theorem]{Lemma}
\newtheorem{problem}[theorem]{Open problem}
\theoremstyle{definition}
\newtheorem{definition}[theorem]{Definition}
\theoremstyle{remark}
\newtheorem{remark}[theorem]{Remark}

\newcommand{\C}{\mathbb C}
\newcommand{\R}{\mathbb R}
\newcommand{\N}{\mathbb N}
\newcommand{\cA}{\mathcal A}
\newcommand{\cH}{\mathcal H}
\newcommand{\cK}{\mathcal K}
\newcommand{\cP}{\mathcal P}
\newcommand{\cQ}{\mathcal Q}
\newcommand{\Rea}{\operatorname{Re}}
\newcommand{\Ima}{\operatorname{Im}}
\newcommand{\Span}{\operatorname{span}}
\newcommand{\dist}{\operatorname{dist}}
\newcommand{\norm}[1]{\lVert #1\rVert}
\newcommand{\ip}[2]{\langle #1,#2\rangle}

\title{A Universal One-Probe Real-Descent Law\\
and the Riemann Hypothesis}
\author{}
\date{Draft of July 22, 2026}

\begin{document}
\maketitle

\begin{abstract}
Let $\mathcal Z\subset\C\setminus\{0\}$ be any conjugation-invariant multiset and let $z\notin
\mathcal Z$ be nonreal.  Put
\[
 K_z(p)(a)=\frac{p(a^{-1})}{a-z},
 \qquad
 r_z(a)=\frac{\bar a^{-2}}{\bar a-z}.
\]
We prove a universal pointwise field-descent law:
\[
 \exists (p_n)\subset\R[X]\ \forall a\in\mathcal Z,
       K_z(p_n)(a)\longrightarrow r_z(a)
 \quad\Longleftrightarrow\quad
 \mathcal Z\subset\R.
\]
No discreteness, summability, norm convergence, or positivity hypothesis is used; if the right
side holds, the constant choice $p_n=X^2$ is exact.  Thus reality of an entire symmetric spectrum
is encoded by descent of one pointwise Cauchy response from complex to real polynomial scalars.

Let $\rho$ run through the nontrivial zeros of the Riemann zeta function, with multiplicity, and
put $a_\rho=-i(\rho-\tfrac12)$.  Thus the Riemann hypothesis says that every $a_\rho$ is real.
For a regular point $z=u+iv$ define the single Cauchy-probe defect
\[
 \mathfrak D(z)=
 \sum_\rho\frac{-\Ima(a_\rho^{-2})}{|a_\rho-z|^2}.
\]
We prove the exact quartet-variance identity
\[
 \mathfrak D(z)=16uv\sum_\rho
 \frac{(\Rea a_\rho)^2(\Ima a_\rho)^2
       (|a_\rho|^2+|z|^2)}
 {|a_\rho|^4|a_\rho-z|^2|\overline{a_\rho}-z|^2
  |a_\rho+z|^2|\overline{a_\rho}+z|^2}.
\]
At the fixed zero-free point $z_0=(1+i)/4$ the prefactor is one.  Consequently
\[
 \mathfrak D(z_0)\geq0,
 \qquad
 \mathfrak D(z_0)=0\Longleftrightarrow\mathrm{RH},
 \qquad
 \mathfrak D(z_0)>0\Longleftrightarrow\neg\mathrm{RH}.
\]
This replaces a family of positivity tests by one predetermined scalar with an unconditional
sign and a cancellation-free failure decomposition.

The same scalar has a one-vector operator interpretation.  On the multiplicity-exact zero space
let $B$ be multiplication by $a_\rho^{-1}$, let
$k_z(a)=(a-z)^{-1}$, and let
$r_z(a)=\bar a^{-2}(\bar a-z)^{-1}$.  Complex polynomials always recover $r_{z_0}$ from the
single Krylov orbit of $k_{z_0}$, whereas real polynomials do so exactly under RH:
\[
 r_{z_0}\in\overline{\{p(B)k_{z_0}:p\in\R[X]\}}
 \quad\Longleftrightarrow\quad\mathrm{RH}.
\]
The obstruction is local and explicit: every real-Krylov limit satisfies a two-point seam law,
and its residual on $r_{z_0}$ is a nonzero rational factor times $\Ima a_\rho$.
Finally, we connect the criterion to finite divided differences of the continued
prime--Gamma logarithmic derivative.  This isolates a new source-side route: construct one real
reciprocal-Krylov recovery, or prove the opposite sign of one fixed prime--Gamma energy contrast.
An attempted canonical realification of the unconditional complex recovery is computed exactly
and shown to stop at the same seam, so conjugation symmetry alone cannot close the route.  We do
not prove the remaining source theorem and do not claim a proof of RH.
\end{abstract}

\paragraph{Logical status.}
Every theorem, proposition, lemma, and corollary below is unconditional.  The only unproved
statement is explicitly labelled as an open problem.  In particular, the paper proves a new
one-probe criterion and an exact real-form obstruction; it does not prove the Riemann hypothesis.

\section{The one-probe principle}

Write
\[
 \xi(s)=\frac12s(s-1)\pi^{-s/2}\Gamma(s/2)\zeta(s),
 \qquad
 X(z)=\xi\!\left(\frac12+iz\right).
\]
The zeros of $X$ form a discrete multiset $\cA$, counted with analytic multiplicity.  We use
letters $a,x+iy\in\cA$.  The functional equation and Schwarz symmetry give multiplicity-preserving
involutions
\begin{equation}
 a\longmapsto-a,
 \qquad
 a\longmapsto\bar a.                                      \tag{1.1}
\end{equation}
No point of $\cA$ lies on the imaginary axis: in the original $s$-coordinate this is the
classical fact that $\zeta$ has no real zero in $0<s<1$.  Moreover
\begin{equation}
 \sum_{a\in\cA}|a|^{-2}<\infty.                            \tag{1.2}
\end{equation}
The standard symmetric enumeration is not needed below because every series that is rearranged
will be absolutely convergent.

For $z\notin\cA$, define
\begin{equation}
 d_z(a)=-\frac{\Ima(a^{-2})}{|a-z|^2},
 \qquad
 \mathfrak D(z)=\sum_{a\in\cA}d_z(a).                    \tag{1.3}
\end{equation}
For large $|a|$, the summand is $O_z(|a|^{-4})$; discreteness handles the finite inner part.
Thus (1.3) converges absolutely.

The key algebra is a four-point identity.

\begin{lemma}[Quartet variance]
\label{lem:quartet}
Let $a=x+iy\ne0$ and $z=u+iv$, and suppose none of
$a,\bar a,-a,-\bar a$ equals $z$.  Then
\begin{align}
 &d_z(a)+d_z(\bar a)+d_z(-a)+d_z(-\bar a) \notag\\
 &\quad=64uv\,
 \frac{x^2y^2(|a|^2+|z|^2)}
 {|a|^4|a-z|^2|\bar a-z|^2|a+z|^2|\bar a+z|^2}.          \tag{1.4}
\end{align}
\end{lemma}

\begin{proof}
Since
\[
 -\Ima(a^{-2})=\frac{2xy}{|a|^4},
\]
the left side of (1.4) equals
\[
 \frac{2xy}{|a|^4}
 \left(
 \frac1{|a-z|^2}-\frac1{|\bar a-z|^2}
 +\frac1{|a+z|^2}-\frac1{|\bar a+z|^2}
 \right).
\]
Putting the four fractions over their common denominator, the numerator is
\[
 32uvxy(|a|^2+|z|^2).
\]
Multiplication by the outer factor $2xy/|a|^4$ gives (1.4).
\end{proof}

The calculation is not specific to zeta.  It gives the following universal law.

\begin{theorem}[Universal one-probe law]
\label{thm:universal}
Let $\mathcal Z$ be a discrete multiset in $\C\setminus\{0\}$, invariant with multiplicity under
$a\mapsto-a$ and $a\mapsto\bar a$.  Assume that the series below are absolutely convergent and
that $z=u+iv$ misses $\mathcal Z$.  Then
\begin{equation}
 \sum_{a\in\mathcal Z}\frac{-\Ima(a^{-2})}{|a-z|^2}
 =16uv\sum_{a\in\mathcal Z}
 \frac{(\Rea a)^2(\Ima a)^2(|a|^2+|z|^2)}
 {|a|^4|a-z|^2|\bar a-z|^2|a+z|^2|\bar a+z|^2}.          \tag{UQV}
\end{equation}
If $uv\ne0$ and $\mathcal Z$ avoids the imaginary axis, the left side vanishes exactly when
$\mathcal Z\subset\R$.
\end{theorem}

\begin{proof}
Sum Lemma~\ref{lem:quartet} over $\mathcal Z$.  Each of its four left-hand sums is the same by
the two multiplicity-preserving involutions, so division by four gives (UQV).  Every summand on
the right is nonnegative.  Under the axis-avoidance hypothesis, its numerator vanishes exactly
when $\Ima a=0$.
\end{proof}

Define the nonnegative quartet variance
\begin{equation}
 V_z(a)=
 \frac{(\Rea a)^2(\Ima a)^2(|a|^2+|z|^2)}
 {|a|^4|a-z|^2|\bar a-z|^2|a+z|^2|\bar a+z|^2}.          \tag{1.5}
\end{equation}

\begin{theorem}[One-probe quartet-variance law]
\label{thm:variance}
For every $z=u+iv\notin\cA$,
\begin{equation}
 \boxed{\quad \mathfrak D(z)=16uv\sum_{a\in\cA}V_z(a).\quad} \tag{1.6}
\end{equation}
In particular,
\begin{equation}
 uv\,\mathfrak D(z)\geq0.                                 \tag{1.7}
\end{equation}
If $uv\ne0$, then
\begin{equation}
 \mathfrak D(z)=0
 \quad\Longleftrightarrow\quad
 \Ima a=0\text{ for every }a\in\cA
 \quad\Longleftrightarrow\quad\mathrm{RH}.              \tag{1.8}
\end{equation}
\end{theorem}

\begin{proof}
Absolute convergence permits reindexing by both involutions in (1.1).  Hence the sum over
$\cA$ of the left side of (1.4) is $4\mathfrak D(z)$, while its right side is
$64uv\sum_aV_z(a)$.  This proves (1.6), and (1.7) follows termwise.

Assume $uv\ne0$.  Every denominator in (1.5) is positive.  Also $\Rea a\ne0$ for all
$a\in\cA$.  Thus $V_z(a)=0$ exactly when $\Ima a=0$.  A convergent sum of nonnegative terms
vanishes exactly when every term vanishes.  Finally,
$a=-i(\rho-\tfrac12)$ is real exactly when $\Rea\rho=1/2$.
\end{proof}

The point used throughout the paper is
\begin{equation}
 z_0=\frac{1+i}{4}.                                      \tag{1.9}
\end{equation}
It is regular because $X$ has no zero at the corresponding low height.  More importantly,
\[
 16\Rea(z_0)\Ima(z_0)=1.
\]

\begin{corollary}[Predetermined scalar witness]
\label{cor:fixed}
At $z_0=(1+i)/4$,
\begin{equation}
 \boxed{
 \mathfrak D(z_0)=\sum_{a\in\cA}V_{z_0}(a)\geq0,
 \qquad
 \mathfrak D(z_0)=0\Longleftrightarrow\mathrm{RH}.}       \tag{1.10}
\end{equation}
Equivalently,
\begin{equation}
 \mathfrak D(z_0)>0\Longleftrightarrow\neg\mathrm{RH},
 \qquad
 \mathfrak D(z_0)\leq0\Longleftrightarrow\mathrm{RH}.    \tag{1.11}
\end{equation}
\end{corollary}

The unusual feature of (1.10) is not merely that one scalar is equivalent to RH.  Its sign is
known in advance, and every hypothetical off-line quartet contributes with the same strict sign.
There is no cancellation among distinct heights, no search for a witnessing test function, and
no limit over a family of inequalities.

\section{Operator form: one vector, not one operator}

Let $\cH=\ell^2(\cA)$, with multiplicities retained, and take the inner product conjugate-linear
in the first variable.  Define compact diagonal operators
\begin{equation}
 (Bf)(a)=a^{-1}f(a),
 \qquad
 (Af)(a)=a^{-2}f(a)=B^2f(a).                              \tag{2.1}
\end{equation}
The reciprocal summability (1.2) makes $B$ Hilbert--Schmidt and $A$ trace class.  For a regular
point $z$, set
\begin{equation}
 k_z(a)=\frac1{a-z}.                                      \tag{2.2}
\end{equation}
Then $k_z\in\cH$ and
\begin{equation}
 \mathfrak D(z)=\Ima\ip{Ak_z}{k_z}.                      \tag{2.3}
\end{equation}
Thus Corollary~\ref{cor:fixed} is a one-vector weak-symmetry criterion.  It does not require
$A$ to be self-adjoint on all of $\cH$; reality of one predetermined quadratic value suffices.

Let $J$ be the complex-linear involution induced by $a\mapsto\bar a$:
\begin{equation}
 (Jf)(a)=f(\bar a).                                       \tag{2.4}
\end{equation}
It is a unitary involution.  Define the adjoint response
\begin{equation}
 r_z=JAk_z,
 \qquad
 r_z(a)=\frac{\bar a^{-2}}{\bar a-z}.                    \tag{2.5}
\end{equation}

\begin{proposition}[Polarization of the fixed defect]
\label{prop:polarization}
Put $Q(f)=\Rea\ip{Jf}{f}$.  Then
\begin{equation}
 \frac{Q(k_z+i r_z)-Q(k_z-i r_z)}4=\mathfrak D(z).        \tag{2.6}
\end{equation}
\end{proposition}

\begin{proof}
Expand both quadratic values.  Their diagonal terms cancel, and $Jr_z=Ak_z$ gives the remaining
cross term $\Ima\ip{Ak_z}{k_z}$.
\end{proof}

Formula (2.6) is the bridge to the Weil form: $Q$ is the zero-side Hermitian form obtained by
critical-line reflection.  The next section shows that the decisive distinction is not existence
of a complex cyclic realization, but descent of that realization to real coefficients.

\section{Universal complex cyclicity}

For $p\in\C[X]$ define
\begin{equation}
 K_z(p)=p(B)k_z,
 \qquad
 K_z(p)(a)=p(a^{-1})(a-z)^{-1}.                           \tag{3.1}
\end{equation}
Vectors coming from functions of the spectral location are constant on all multiplicity copies
of a fixed zero.  Call their closed subspace $\cH_{\rm ev}$.

\begin{theorem}[One-Cauchy-vector complex cyclicity]
\label{thm:cyclic}
For every regular $z$,
\begin{equation}
 \overline{\{K_z(p):p\in\C[X]\}}=\cH_{\rm ev}.           \tag{3.2}
\end{equation}
In particular, there are complex polynomials $p_n$ such that
\begin{equation}
 K_{z_0}(p_n)\longrightarrow r_{z_0}\quad\text{in }\cH. \tag{3.3}
\end{equation}
No zero-location hypothesis is used.
\end{theorem}

\begin{proof}
Write $\mu_a=a^{-1}$.  For every $\delta>0$, the outer spectral shell
$\{a:|\mu_a|\geq\delta\}$ is finite.  Fix one spectral location $\alpha=\mu_a$ and a desired
fiber value $c$.  On a finite outer shell, a Lagrange polynomial separates $\alpha$ from every
other spectral location.  Multiplying by a high power of $X/\alpha$ leaves its value at
$\alpha$ unchanged and drives its value to zero at every point with modulus strictly below
$|\alpha|$.  Enlarging the shell if necessary handles ties of equal modulus.  The resulting
polynomials are uniformly bounded on the reciprocal spectrum.  Dominated convergence against
$|k_z|^2$ therefore recovers the complete multiplicity fiber at $a$ in norm.

Finite sums recover all finitely supported fiber-constant vectors, and such vectors are dense in
$\cH_{\rm ev}$.  This proves (3.2).  The response (2.5) is square-summable and fiber-constant,
so (3.3) follows.
\end{proof}

Theorem~\ref{thm:cyclic} is an important guardrail.  A complex Krylov realization of the adjoint
response cannot prove RH: it exists unconditionally.  The arithmetic content starts only after
imposing the real form.

\section{Real Krylov descent is exactly RH}

Every real polynomial obeys $p(\bar w)=\overline{p(w)}$.  Applied at a conjugate pair of spectral
locations, this gives an exact seam equation.

\begin{lemma}[Real-Krylov seam]
\label{lem:seam}
If $p\in\R[X]$ and $v=K_z(p)$, then for every $a\in\cA$,
\begin{equation}
 (\bar a-z)v(\bar a)
 -(\bar a-\bar z)\overline{v(a)}=0.                       \tag{4.1}
\end{equation}
The same identity holds for every norm limit of real-Krylov vectors.
\end{lemma}

\begin{proof}
For a polynomial vector,
\[
 (\bar a-z)v(\bar a)=p(\bar a^{-1})
 =\overline{p(a^{-1})}
 =(\bar a-\bar z)\overline{v(a)}.
\]
Coordinate evaluation is continuous on $\ell^2(\cA)$, so (4.1) survives norm limits.
\end{proof}

The seam residual of the desired response has no hidden term.

\begin{lemma}[Adjoint-response seam residual]
\label{lem:residual}
For $r_z$ from (2.5), the left side of (4.1) is
\begin{equation}
 \frac{4\,\Ima(a)\Ima(z)\,a^{-2}}
 {(a-z)(a-\bar z)}.                                      \tag{4.2}
\end{equation}
\end{lemma}

\begin{proof}
Substitute
\[
 r_z(\bar a)=\frac{a^{-2}}{a-z},
 \qquad
 \overline{r_z(a)}=\frac{a^{-2}}{a-\bar z}
\]
into (4.1), put the two terms over a common denominator, and use
$(\bar a-a)(\bar z-z)=4\Ima(a)\Ima(z)$.
\end{proof}

The seam is not special to the zeta divisor, to a Hilbert norm, or even to a discrete spectrum.
It gives the following field-descent classification.

\Needspace{17\baselineskip}
\begin{theorem}[Universal pointwise real-descent law]
\label{thm:universaldescent}
Let $\mathcal Z\subset\C\setminus\{0\}$ be a multiset invariant with multiplicity under complex
conjugation.  Let $z\notin\mathcal Z$ with $\Ima z\ne0$.  For
\[
 K_z(p)(a)=\frac{p(a^{-1})}{a-z},
 \qquad
 r_z(a)=\frac{\bar a^{-2}}{\bar a-z},
\]
the following are equivalent:
\begin{enumerate}
 \item there is one sequence $p_n\in\R[X]$ such that
       $K_z(p_n)(a)\to r_z(a)$ for every $a\in\mathcal Z$;
 \item $\mathcal Z\subset\R$.
\end{enumerate}
The forward implication requires only pointwise convergence.  Under the second condition the
constant sequence $p_n=X^2$ gives equality at every point.
\end{theorem}

\begin{proof}
Every $K_z(p_n)$ satisfies the seam (4.1).  Pointwise convergence at $a$ and $\bar a$ therefore
passes that identity to $r_z$.  Lemma~\ref{lem:residual} gives
\[
 \frac{4\,\Ima(a)\Ima(z)\,a^{-2}}
 {(a-z)(a-\bar z)}=0.
\]
All factors except $\Ima(a)$ are nonzero, so $\Ima(a)=0$.  This holds for every
$a\in\mathcal Z$.  Conversely, if $a=\bar a$, then
$r_z(a)=a^{-2}(a-z)^{-1}=K_z(X^2)(a)$.
\end{proof}

\begin{theorem}[One-vector real-descent law]
\label{thm:realdescent}
At the fixed point $z_0=(1+i)/4$,
\begin{equation}
 \boxed{\quad
 r_{z_0}\in\overline{\{p(B)k_{z_0}:p\in\R[X]\}}
 \quad\Longleftrightarrow\quad\mathrm{RH}.
 \quad}                                                   \tag{4.3}
\end{equation}
More precisely, complex polynomial recovery always holds by
Theorem~\ref{thm:cyclic}, while real polynomial recovery holds under RH with the single polynomial
$p(X)=X^2$.
\end{theorem}

\begin{proof}
Suppose real polynomial vectors converge to $r_{z_0}$.  Lemma~\ref{lem:seam} gives (4.1) for the
limit.  In (4.2), $\Ima z_0\ne0$, and all remaining rational factors are nonzero at a regular
probe.  Hence $\Ima a=0$ for every $a\in\cA$, which is RH.

Conversely, under RH every $a$ is real.  Then $J$ acts trivially on spectral locations, and
\[
 r_{z_0}=Ak_{z_0}=B^2k_{z_0}=K_{z_0}(X^2).
\]
\end{proof}

Theorem~\ref{thm:realdescent} is stronger than the statement that a hypothetical Hilbert--P\'olya
operator should be self-adjoint.  The compact diagonal operator, its spectrum, a cyclic vector,
and the target response already exist without RH.  The entire hypothesis is concentrated in one
descent question:
\[
 \textit{does one known complex Krylov solution descend to the real coefficient field?}
\]

\section{Prime--Gamma source realization}

The preceding sections use the divisor to expose the geometry.  A route toward RH must recover
the decisive real structure from the arithmetic source.  We now state the exact interface.

\Needspace{12\baselineskip}
Let
\[
 s(z)=\frac12+iz,
 \qquad
 A_\infty(s)=\frac12s(s-1)\pi^{-s/2}\Gamma(s/2),
\]
and define the meromorphically continued prime--Gamma trace
\begin{equation}
 T(z)=-i\left(-\frac{\zeta'}{\zeta}(s(z))
                   -\frac{A_\infty'}{A_\infty}(s(z))\right).
                                                               \tag{5.1}
\end{equation}
At every regular point this is $i\xi'(s(z))/\xi(s(z))$, equivalently $X'(z)/X(z)$.
For regular, conjugate-separated points define
\begin{equation}
 \cK(z,w)=\frac{T(\bar w)-T(z)}{z-\bar w}.                \tag{5.2}
\end{equation}
In the absolute Euler half-plane, (5.1) consists of the convergent von Mangoldt series and the
explicit Gamma connection.  Formula (5.2) then continues this source kernel into the critical
strip.

For a finite packet $P=\{(c_j,z_j)\}_{j=1}^m$, put
\begin{equation}
 E(P)=\Rea\sum_{i,j=1}^m c_i\bar c_j\cK(z_i,z_j).         \tag{5.3}
\end{equation}

\begin{proposition}[Finite source identity]
\label{prop:source}
Let $F_P(a)=\sum_jc_j(a-z_j)^{-1}$.  Whenever all denominators in (5.2) are regular,
\begin{equation}
 E(P)=\Rea\sum_{a\in\cA}F_P(a)\overline{F_P(\bar a)}
     =Q(F_P|_{\cA}).                                     \tag{5.4}
\end{equation}
\end{proposition}

\begin{proof}
The Hadamard logarithmic derivative of $X$ gives the divided-difference identity
\[
 \cK(z,w)=\sum_{a\in\cA}\frac{1}{(a-z)(a-\bar w)}.
\]
Absolute convergence permits expansion of the finite double sum in (5.3), yielding (5.4).
The equality with the prime--Gamma expression follows from the completed logarithmic derivative
in (5.1).
\end{proof}

Combining Cauchy density, Proposition~\ref{prop:source}, and continuity of $Q$ gives an arithmetic
approximation of the one-probe defect.

\begin{theorem}[Finite prime--Gamma contrast reconstruction]
\label{thm:reconstruction}
There are two sequences of finite, conjugate-separated regular packets $P_n^+$ and $P_n^-$ such
that their spectral evaluations converge respectively to $k_{z_0}+ir_{z_0}$ and
$k_{z_0}-ir_{z_0}$, and
\begin{equation}
 \frac{E(P_n^+)-E(P_n^-)}4\longrightarrow\mathfrak D(z_0). \tag{5.5}
\end{equation}
Consequently, eventual nonpositivity of the contrast in (5.5) implies RH.
\end{theorem}

\begin{proof}
Finite regular Cauchy packets are dense in the fiber-constant evaluation space.  Colliding poles
may be perturbed to conjugate-separated regular poles without changing the Hilbert limit.
Apply this to the two square-summable target vectors.  Proposition~\ref{prop:source} identifies
each finite energy with $Q$, and continuity together with (2.6) gives (5.5).  Corollary
\ref{cor:fixed} supplies the final implication.
\end{proof}

The packets in Theorem~\ref{thm:reconstruction} are existence witnesses extracted from spectral
closure.  That fact matters: the theorem does not yet provide a divisor-free recipe for their
coefficients.  It does, however, prove that no analytic continuation, pole separation, or finite
source representation gap remains after such a recipe is found.

\Needspace{12\baselineskip}
\section{The new route and its exact open core}

The results give two sufficient ways to attack the remaining source problem.

\begin{problem}[Real source selection]
\label{prob:real}
Construct, from the continued prime--Gamma trace $T$ and its divided differences alone, real
polynomials $p_n\in\R[X]$ such that
\begin{equation}
 p_n(B)k_{z_0}\longrightarrow r_{z_0}.                   \tag{6.1}
\end{equation}
Alternatively, construct source-determined packets in Theorem~\ref{thm:reconstruction} whose
energy contrast is eventually nonpositive.
\end{problem}

By Theorem~\ref{thm:realdescent}, a solution of Problem~\ref{prob:real} proves RH.  The point is
not that the problem is known to be easier.  Its value is that it exposes a sharply different
mechanism from an all-test positivity theorem:
\begin{equation}
 \boxed{
 \begin{gathered}
 \text{one fixed probe}\longrightarrow\text{one adjoint response}\\
 \longrightarrow\text{real descent of one complex Krylov orbit}
 \longrightarrow\mathrm{RH}.
 \end{gathered}}                                         \tag{6.2}
\end{equation}

The source data needed by an iterative construction are explicit divided differences of (5.1).
The missing invariant is coefficient reality.  This suggests studying conjugation-preserving
continued-fraction, Lanczos, or moment recursions for the prime--Gamma kernel, but any such
recursion must prove its realness rather than import self-adjointness.

\section{Hostile audit}

We record what the new formulation does and does not accomplish.

\subsection*{Complex cyclicity is not evidence for RH}
The complex orbit is universal by Theorem~\ref{thm:cyclic}.  Any argument ending with
$r_{z_0}\in\overline{\C[B]k_{z_0}}$ has proved no zero-location information.  Complex Lagrange
coefficients absorb every off-line phase.

\subsection*{Finite real interpolation cannot evade the wall}
At a conjugate pair $a,\bar a$, real coefficients force (4.1).  The target residual (4.2) is
already nonzero when $\Ima a\ne0$.  Therefore neither higher polynomial degree nor a more
accurate finite interpolation can solve the real problem under a hypothetical off-line zero.
The obstruction is closed under norm limits.

\subsection*{Canonical realification was attempted and fails exactly at the frontier}
Start with the unconditional complex recovery $K_{z_0}(p_n)\to r_{z_0}$.  Let $p_n^\#$ be
obtained by conjugating every coefficient and let
$q_n=(p_n+p_n^\#)/2\in\R[X]$.  Direct coefficient transport gives, at every regular coordinate,
\begin{equation}
 K_z(q_n)(a)=\frac12\left(
 K_z(p_n)(a)+\frac{a-\bar z}{a-z}\,
       \overline{K_z(p_n)(\bar a)}\right).               \tag{7.1}
\end{equation}
Hence $K_z(q_n)$ converges pointwise to the canonical real-form projection
\[
 (P_zv)(a)=\frac12\left(v(a)+\frac{a-\bar z}{a-z}\overline{v(\bar a)}\right)
\]
with $v=r_z$.  The equality $P_zr_z=r_z$ is exactly the seam equation (4.1), and by (4.2) it
holds at every coordinate exactly when all coordinates are real.  Thus the most direct attempt
to turn the known complex source solution into a real one succeeds exactly under RH and cannot
be used as an unconditional proof.  This calculation rules out coefficient averaging as the
missing source mechanism; any successful construction must inject additional zeta-specific
arithmetic structure before realification.

\subsection*{The known sign points in the wrong direction}
Equation (1.10) proves $\mathfrak D(z_0)\geq0$ unconditionally.  More estimates of the same sign
cannot prove RH.  A source argument must force the opposite inequality, exact cancellation, or
real Krylov descent.  This prevents a positive-resolvent or Cauchy--Schwarz argument from being
mistaken for progress.

\subsection*{Weil Hermiticity is insufficient}
The form $Q(f)=\Rea\ip{Jf}{f}$ is Hermitian because $J$ is a unitary involution.  Hermiticity does
not identify $J$ with the identity and does not make $A$ self-adjoint in the ordinary Hilbert
metric.  The seam (4.1) measures exactly that missing identification.

\subsection*{The source reconstruction is not yet canonical}
Theorem~\ref{thm:reconstruction} removes a representation gap but selects approximants through
the target Hilbert closure.  Treating their existence as a source proof would be circular.  A
valid proof must specify coefficients from prime--Gamma data and then establish the real-form
constraint independently.

These failures do not invalidate route (6.2).  They isolate its only live producer: a
conjugation-preserving source selection theorem.  Conversely, they make the route falsifiable.
A claimed construction can be rejected immediately if it uses complex coefficients, a
zero-selected interpolation set, ordinary Hilbert self-adjointness, or the desired sign as an
input.

\section{Relation to existing criteria}

Li's criterion replaces RH by positivity of an infinite sequence of coefficients, and
Bombieri--Lagarias relate that sequence to Weil positivity~\cite{Li1997,BombieriLagarias1999}.
Weil's criterion asks for nonnegativity on a full test class~\cite{Weil1952}.  Recent spectral
programs seek self-adjoint realizations or limits of self-adjoint finite systems; see, for example,
Hedenmalm's theta differential pair~\cite{Hedenmalm2026} and Suzuki's screw-function framework
\cite{Suzuki2026}.  Caruso describes closed Krylov spaces of compact normal operators through
their spectral projections~\cite{Caruso2023}, while general Krylov solvability is closely related
to spectral measures and moment problems~\cite{Caruso2024}.

The present criterion has a different quantifier structure.  It fixes one regular point, one
Cauchy vector, one compact reciprocal operator, and one target response.  RH is neither positivity
of all vectors nor existence of a complex cyclic model.  It is descent of one already-existing
complex solution from $\C[X]$ to $\R[X]$.  The universal pointwise theorem shows that this field
descent, unlike ordinary compact-normal Krylov solvability, is already rigid before a Hilbert
topology is imposed.  The quartet identity then supplies a cancellation-free one-scalar witness
for failure.

\enlargethispage{\baselineskip}
Targeted searches across the Li/Weil, Cauchy-transform, Hilbert--P\'olya, and Krylov literature
found no prior statement of the universal pointwise descent theorem, the fixed-probe identity
(1.10), or the zeta real-descent equivalence (4.3).  This is not a claim of exhaustive priority.
The algebraic ingredients are elementary; the claimed contribution is their field-sensitive,
multiplicity-exact synthesis and the resulting one-vector source interface.

\section{Conclusion}

The central result is the universal classification in
Theorem~\ref{thm:universaldescent}: one nonreal Cauchy response descends pointwise from
$\C[X]$ to $\R[X]$ exactly when the entire conjugation-symmetric spectrum is real.  For the
centered zeta divisor it yields the paired fixed-probe law
\[
 \boxed{
 \begin{gathered}
 \mathfrak D\!\left(\frac{1+i}{4}\right)
   =\sum_{a\in\cA}V_{(1+i)/4}(a)\geq0,\\
 \mathfrak D\!\left(\frac{1+i}{4}\right)=0
   \Longleftrightarrow
 r_{(1+i)/4}\in\overline{\R[B]k_{(1+i)/4}}
   \Longleftrightarrow\mathrm{RH}.
 \end{gathered}}                                         \tag{9.1}
\]
Every off-line quartet contributes strictly to the first scalar and creates an explicit local
seam residual in the second criterion.  Yet the same response always lies in the complex Krylov
closure.  RH is therefore the exact real-form descent of one universal complex cyclic response.

This does not solve RH.  It creates a new route with one unambiguous frontier: obtain coefficient
reality from the continued prime--Gamma divided-difference source.  Success would close the chain;
failure can no longer hide behind analytic continuation, density, finite packet realization,
operator existence, or a family of unlocated negative directions.

\begin{thebibliography}{99}

\bibitem{BombieriLagarias1999}
E.~Bombieri and J.~C.~Lagarias,
\emph{Complements to Li's criterion for the Riemann hypothesis},
Journal of Number Theory \textbf{77} (1999), 274--287,
\href{https://doi.org/10.1006/jnth.1999.2392}{doi:10.1006/jnth.1999.2392}.

\bibitem{Caruso2024}
N.~A.~Caruso,
\emph{Structural properties of Krylov subspaces, Krylov solvability, and applications to
unbounded self-adjoint operators},
arXiv:2404.19698 (2024),
\url{https://arxiv.org/abs/2404.19698}.

\bibitem{Caruso2023}
N.~A.~Caruso,
\emph{A note on the Krylov solvability of compact normal operators on Hilbert space},
Complex Analysis and Operator Theory \textbf{17} (2023), Article 100,
\url{https://doi.org/10.1007/s11785-023-01413-0}.

\bibitem{Hedenmalm2026}
H.~Hedenmalm,
\emph{Spectral interpretation of Riemann zeta zeros},
arXiv:2606.17494 (2026),
\url{https://arxiv.org/abs/2606.17494}.

\bibitem{Li1997}
X.-J.~Li,
\emph{The positivity of a sequence of numbers and the Riemann hypothesis},
Journal of Number Theory \textbf{65} (1997), 325--333.

\bibitem{Suzuki2026}
M.~Suzuki,
\emph{Weil's quadratic form via the screw function},
arXiv:2606.09096 (2026),
\url{https://arxiv.org/abs/2606.09096}.

\bibitem{Titchmarsh1986}
E.~C.~Titchmarsh, revised by D.~R.~Heath-Brown,
\emph{The Theory of the Riemann Zeta-Function}, second edition,
Oxford University Press, 1986.

\bibitem{Weil1952}
A.~Weil,
\emph{Sur les formules explicites de la th\'eorie des nombres premiers},
Communications du S\'eminaire Math\'ematique de l'Universit\'e de Lund (1952), 252--265.

\end{thebibliography}

\end{document}
